% ======================================================================
% Capítulo: Metodología y planificación
% Scrum individual + SDD con OpenSpec, calidad y planificación real frente a la del anteproyecto. Extensión orientativa: 3--4 páginas.
% ======================================================================
\chapter{Metodología y planificación}
\label{cap:metodologia}

\section{Scrum adaptado a un desarrollo individual}
\label{sec:scrum}
\label{sec:iteraciones}

El trabajo se ha organizado con Scrum \cite{schwaber2020}, adaptado a un desarrollo individual: el autor es a la vez propietario del producto y equipo de desarrollo, y los tutores, las partes interesadas. Cada \emph{sprint}, de varias semanas, parte de un objetivo que se traduce en cambios de OpenSpec (sección~\ref{sec:sdd}), cuyas tareas hacen de pila del \emph{sprint}, y se revisa comprobando el objetivo en la aplicación desplegada. Tras el contacto con el Ayuntamiento de Sedella (sección~\ref{sec:motivacion}), la prioridad fue un producto mínimo viable (MVP) que enseñar a entidades públicas, así que la integración continua y el despliegue se prepararon en julio y agosto de 2026, en lugar de al final. A principios de septiembre la primera versión se abrió durante dos días a algunas personas, de manera informal y sin protocolo, y sus comentarios se centraron en el tiempo de respuesta y en la usabilidad; tras conocer el coste de la primera ingesta en la nube (sección~\ref{sec:recursos}) se aplazó la apertura a usuarios, y en noviembre se renunció a ella dentro de este trabajo (sección~\ref{sec:validacion-usuarios}). Los \emph{sprints} siguientes ampliaron y consolidaron el sistema (sección~\ref{sec:desarrollo-real}).

\section{Desarrollo guiado por especificaciones}
\label{sec:sdd}

El desarrollo guiado por especificaciones (\emph{Spec-Driven Development}, SDD) acuerda por escrito qué debe hacer el sistema antes de escribir código y mantiene ese texto como la referencia contra la que se implementa y se revisa. Cuando parte del código lo proponen agentes de IA, la especificación es además el contexto que reciben y el criterio con el que se comprueba su trabajo. Es el planteamiento de OpenSpec \cite{openspec}, la herramienta de código abierto que usa el proyecto, y de otras herramientas como Spec Kit, de GitHub \cite{delimarsky2025speckit}; los trabajos académicos sobre el tema son todavía preprints recientes \cite{piskala2026sdd}.

El directorio \texttt{openspec/} contiene \texttt{config.yaml}, con el contexto del proyecto y las reglas de los artefactos; \texttt{specs/}, con una especificación por \emph{capacidad} del sistema, cuyos requisitos se redactan con \texttt{SHALL} y se ilustran con escenarios \texttt{WHEN}/\texttt{THEN} (ejemplo y lista completa en el apéndice~\ref{ap:especificaciones}), y \texttt{changes/}, con los cambios en curso y los archivados.

Cada cambio sigue el mismo ciclo, que OpenSpec expone como comandos para los asistentes de IA: \textbf{explorar} el problema sin escribir código; \textbf{proponer} el cambio con su propuesta, su diseño, sus tareas y, si altera el comportamiento, especificaciones \emph{delta} con los requisitos añadidos, modificados o eliminados; \textbf{aplicar} las tareas, actualizando los artefactos si aparece algo no previsto; \textbf{sincronizar} los \emph{delta} con las especificaciones principales, y \textbf{archivar} el cambio con su fecha.

\subsection{Validación humana y especificaciones vivas}
\label{sec:sdd-validacion}
\label{sec:sdd-vivas}

La IA no trabaja por su cuenta: cada artefacto que genera pasa por la validación del autor antes del siguiente paso. Las especificaciones que redacta el asistente son un borrador que el autor revisa y corrige antes de implementar. El código se entrega siempre en una \emph{pull request}, que revisa primero un agente revisor y después el autor, con una revisión manual del código y una prueba manual de las acciones afectadas, además de la integración continua. Las especificaciones se tratan como documentación viva: cada cambio de comportamiento las actualiza en la misma \emph{pull request}, y una diferencia con el código se corrige en ese momento, de modo que \texttt{openspec/specs} describe siempre el sistema completo.

\section{Gestión del código y calidad}
\label{sec:calidad}

Cada cambio se desarrolla en una rama corta y entra en la principal, que es siempre la versión desplegada, mediante una \emph{pull request} (flujo \emph{trunk-based}). Antes de fusionar debe pasar la integración continua (\texttt{.github/workflows/ci.yml}), que en unos dos minutos comprueba la higiene del repositorio, el formato, la compilación y las pruebas de .NET, Angular y Go (con \texttt{go vet}, \texttt{gofmt} y el detector de condiciones de carrera) y que el código generado a partir de los \texttt{.proto} esté al día. El despliegue (\texttt{deploy.yml}) solo se lanza cuando la integración continua pasa en la rama principal. El código que proponen los agentes de IA pasa por los mismos filtros que el escrito a mano.

\section{Uso de la inteligencia artificial en el desarrollo}
\label{sec:ia-desarrollo}

Además de los modelos que forman parte del producto (capítulos~\ref{cap:diseno} y~\ref{cap:implementacion}), se han usado dos asistentes de IA como herramienta de desarrollo. Gemini \cite{gemini} y NotebookLM sirvieron como buscador y apoyo para la investigación: localizar y resumir documentación de las fuentes y las tecnologías y resolver dudas puntuales; lo obtenido se trató como punto de partida y se contrastó con la fuente original. Hasta la incorporación de Claude, Gemini CLI y el entorno Antigravity se usaron también para redactar artefactos de OpenSpec y código, con la misma revisión.
Desde octubre de 2026, Claude, a través de Claude Code \cite{claudecode}, trabaja sobre el repositorio con el mismo ciclo de OpenSpec: lee las especificaciones, propone o actualiza los cambios, implementa las tareas en una rama y abre una \emph{pull request} que el autor revisa; cada \emph{commit} en el que participa lleva la línea \texttt{Co-Authored-By: Claude}. Participó, entre otros trabajos, en la auditoría de código de octubre de 2026 y su corrección (autorización, aislamiento entre ayuntamientos, credenciales y mensajería), en el cotejo de las especificaciones con el código, en la separación de la base vectorial tras el puerto \texttt{IVectorStore} y en la observabilidad.

\subsection{Uso de la IA en la redacción de esta memoria y control humano}

En la redacción de esta memoria no se ha dejado a la IA libertad para desarrollar el contenido. Claude se ha utilizado como herramienta de transcripción de voz a texto: el autor dicta sus ideas y el asistente las transcribe y les da formato en \LaTeX, para agilizar la escritura. También se ha usado para revisar el estilo y la coherencia del texto ya escrito. El contenido, la estructura y las conclusiones son del autor, que revisa y corrige cada texto antes de incorporarlo.

En todos estos usos el autor mantiene la responsabilidad sobre el resultado: las especificaciones y el código que genera la IA se validan como se describe en la sección~\ref{sec:sdd-validacion}, y los datos que no proceden del propio proyecto se contrastan con la fuente original.

\section{Planificación temporal}
\label{sec:planificacion}

\subsection{Planificación del anteproyecto}

El anteproyecto, firmado el 2 de octubre de 2026, dividió el trabajo en nueve fases con una dedicación total de 296 horas: planificación y requisitos (23~h), diseño de la arquitectura y del modelo vectorial (25), ingesta en Go (30), estudio de fragmentación (30), \emph{backend} .NET y RAG (50), \emph{frontend} (45), pruebas, integración y Docker (30), validación con usuarios (18) y memoria (45). No se llevó un registro de las horas dedicadas, así que la comparación con lo previsto se limita al calendario.

\subsection{Desarrollo real}
\label{sec:desarrollo-real}

El desarrollo empezó a mediados de julio de 2026, sin trabajo previo, y se extiende hasta la entrega de la memoria, prevista para finales de noviembre. Hasta septiembre se construyó y desplegó el MVP (\emph{scraper}, ingesta vectorial, chat, interfaz y despliegue); en septiembre llegaron la \emph{pipeline} jerárquica, la Arena y el contexto del usuario; en octubre, la reorganización de la capa de aplicación, Qdrant, la auditoría del código, la observabilidad y las variantes del estudio, y en noviembre, las ejecuciones del estudio. El calendario por fases y los hitos de cada mes están en la sección~\ref{sec:calendario}.

\subsection{Desviaciones}

Al trabajar por \emph{sprints} con objetivos funcionales completos, las fases de ingesta, \emph{backend}, \emph{frontend} y pruebas avanzaron a la vez, y la integración continua y el despliegue se prepararon en julio y agosto, mucho antes de lo previsto, para tener pronto un MVP que enseñar (sección~\ref{sec:scrum}). Hubo trabajo no previsto: la migración de pgvector a Qdrant en octubre y necesidades que el anteproyecto no recogía, como la Arena, el soporte de varios ayuntamientos, el control de coste y la observabilidad (sección~\ref{sec:motivacion}). El estudio comparativo se dejó para el final, porque necesitaba las \emph{pipelines} funcionando y un entorno local sin coste por uso, y la validación con usuarios no se llevó a cabo.

\section{Recursos}
\label{sec:recursos}

El trabajo lo ha desarrollado un único alumno, supervisado por el tutor y el cotutor, con herramientas gratuitas o de código abierto. El estudio de fragmentación se ejecutó en un ordenador de sobremesa con una GPU RTX 5070 Ti, y la aplicación se ejecuta en producción en dos servidores virtuales propios que no suponen un coste añadido. El principal coste variable son los proveedores de modelos de lenguaje, que se pagan por \emph{token}: la única ingesta hecha en la nube, para la primera versión en producción, costó unos 30~USD y solo cubría dos de las cuatro variantes. Por eso el desarrollo y la evaluación se hicieron con modelos locales y se aplazó la apertura a usuarios. El detalle del hardware, los servicios y el coste está en la sección~\ref{sec:recursos-detalle}.
